\subsection{UNIVERSAL PROBLEM FOR DERIVATIONS; COMMUTATIVE CASE}

Suppose now that A is a commutative K-algebra and E an A-module; E can be considered as an (A, A)-bimodule whose two external laws are identical with the given A-module law. On the other hand the (A, A)-bimodule structure on \(A \otimes_{K} A\) is identical with its \((A \otimes_{K} A)\) -module structure arising from the commutative ring structure on \(A \otimes_{K} A\) , since in this case, for x, y, u, v in A,

\[
x. (u \otimes v). y = (x u) \otimes (v y) = (x u) \otimes (y v) = (x \otimes y) (u \otimes v).
\]

The kernel \(\mathfrak{J}\) of \(m\) is therefore in this case an ideal of the ring \(\mathbf{A} \otimes_{\mathbb{K}} \mathbf{A}\) and, as \(m: \mathbf{A} \otimes_{\mathbb{K}} \mathbf{A} \to \mathbf{A}\) is surjective, \((\mathbf{A} \otimes_{\mathbb{K}} \mathbf{A}) / \mathfrak{J}\) is isomorphic to \(\mathbf{A}\) ; if also \(\mathbf{E}\) is considered as an \((\mathbf{A} \otimes_{\mathbb{K}} \mathbf{A})\) -module by means of \(m\) (in other words the \((\mathbf{A} \otimes_{\mathbb{K}} \mathbf{A})\) -module \(m_*(\mathbf{E})\) ), the \((\mathbf{A}, \mathbf{A})\) -bimodule homomorphisms \(\mathfrak{J} \to \mathbf{E}\) are identified with the \((\mathbf{A} \otimes_{\mathbb{K}} \mathbf{A})\) -module homomorphisms \(\mathfrak{J} \to \mathbf{E}\) ( \(\S 4, no. 3\) ), in other words there is a canonical \(\mathbf{K}\) -module isomorphism.

\[
\operatorname{Hom} _ {(\mathbf {A}, \mathbf {A})} (\mathfrak{J} , \mathrm{E}) \rightarrow \operatorname{Hom} _ {\mathbf {A} \otimes_ {\mathbf {K}} \mathbf {A}} (\mathfrak{J} , \mathrm{E}).
\]

On the other hand, \(\mathfrak{J}\mathbf{E} = \{0\}\) , for the elements \(1 \otimes x - x \otimes 1\) generate \(\mathfrak{J}\) as an \((\mathbf{A} \otimes_{\mathbb{K}} \mathbf{A})\) -module (no. 10, Lemma 1) and, for all \(z \in \mathbf{E}\) ,

\[
(1 \otimes x - x \otimes 1) z = 0
\]

by virtue of the definition of the \((\mathbf{A} \otimes_{\mathbf{K}} \mathbf{A})\) -module structure on E. Since \(\mathfrak{J}\) is contained in the annihilator of the \((\mathbf{A} \otimes_{\mathbf{K}} \mathbf{A})\) -module E and the \(((\mathbf{A} \otimes_{\mathbf{K}} \mathbf{A})/\mathfrak{J})\) -module structure on E is by definition just the initial A-module structure given on E, there is, taking account of the canonical isomorphism of

\[
\mathfrak{J} \otimes_ {\mathbf {K}} ((\mathbf {A} \otimes_ {\mathbf {K}} \mathbf {A}) / \mathfrak{J})
\]

onto \(\mathfrak{J} / \mathfrak{J}^2\) (§ 4, no. 1, Corollary 1 to Proposition 1), a canonical K-module isomorphism

\[
\operatorname{Hom} _ {\mathbf {A} \otimes_ {\mathbf {K}} \mathbf {A}} (\mathfrak{J} , \mathrm{E}) \rightarrow \operatorname{Hom} _ {\mathbf {A}} (\mathfrak{J} / \mathfrak{J}^ {2}, \mathrm{E}).
\]

Taking account of Proposition 17 of no. 10, it is seen that we have proved the following proposition:

PROPOSITION 18. Let A be a commutative K-algebra and \(\mathfrak{J}\) the ideal the kernel of the surjective canonical homomorphism \(m: \mathbf{A} \otimes_{\mathbf{K}} \mathbf{A} \to \mathbf{A}\) , so that A is isomorphic to \((\mathbf{A} \otimes_{\mathbf{K}} \mathbf{A}) / \mathfrak{J}\) and \(\mathfrak{J} / \mathfrak{J}^2\) has canonically an A-module structure. Let \(d_{\mathbf{A} / \mathbf{K}}: \mathbf{A} \to \mathfrak{J} / \mathfrak{J}^2\) be the K-linear mapping which associates with every \(x \in \mathbf{A}\) the class of \(x \otimes 1 - 1 \otimes x\) modulo \(\mathfrak{J}^2\) . The mapping \(d_{\mathbf{A} / \mathbf{K}}\) is a K-derivation and, for every A-module E and every K-derivation D: \(\mathbf{A} \to \mathbf{E}\) , there exists one and only one A-linear mapping \(g: \mathfrak{J} / \mathfrak{J}^2 \to \mathbf{E}\) such that \(\mathrm{D} = g \circ d_{\mathbf{A} / \mathbf{K}}\) .

The A-module \(\mathfrak{J} / \mathfrak{J}^2\) is called the A-module of K-differentials of A and is denoted by \(\Omega_{\mathbf{K}}(\mathbf{A})\) ; for all \(x \in \mathbf{A}\) , \(d_{\mathbf{A} / \mathbf{K}}(x)\) (also denoted by \(dx\) ) is called the differential of \(x\) ; it has been seen (no. 10, lemma 1) that the elements \(d_{\mathbf{A} / \mathbf{K}}(x)\) , for \(x \in \mathbf{A}\) , form a generating system of the A-module \(\Omega_{\mathbf{K}}(\mathbf{A})\) . Proposition 18 shows that the mapping \(g \mapsto g \circ d_{\mathbf{A} / \mathbf{K}}\) is a canonical A-module isomorphism

\[
\phi_ {\mathbf {A}}: \operatorname{Hom} _ {\mathbf {A}} (\Omega_ {\mathbf {K}} (\mathbf {A}), \mathrm{E}) \rightarrow \mathrm{D} _ {\mathbf {K}} (\mathrm{A}, \mathrm{E})
\]

(the A-module structure on \(D_{K}(A, E)\) being defined by Proposition 7 of no. 5).

The ordered pair \((\Omega_{\mathbf{K}}(\mathbf{A}), d_{\mathbf{A}/\mathbf{K}})\) is therefore the solution of the universal mapping problem where \(\Sigma\) is the species of A-module structure and the \(\alpha\) -mappings the K-derivations from A to an A-module (Set Theory, IV, §3, no. 1).

Example. Let M be a K-module; it follows from Proposition 14 of no. 9 that for every S(M)-module E, the mapping D ↦ D \textbar{} M defines an S(M)-module isomorphism of D\_K(S(M), E) onto Hom\_K(M, E); on the other hand, since E is an S(M)-module, Hom\_K(M, E) is canonically isomorphic to

\[
\operatorname{Hom} _ {\mathbf {S} (\mathbf {M})} (\mathbf {M} \otimes_ {\mathbf {K}} \mathbf {S} (\mathbf {M}), \mathbf {E}),
\]

every K-homomorphism of M into E being uniquely expressible in the form \(x \mapsto h(x \otimes 1)\) , where

\[
h \in \operatorname{Hom} _ {\mathbf {S} (\mathbf {M})} (\mathbf {M} \otimes_ {\mathbf {K}} \mathbf {S} (\mathbf {M}), \mathbf {E})
\]

(II, § 5, no. 1). Let \(\mathbf{D}_0\) be the K-derivation \(\mathbb{S}(\mathbf{M}) \to \mathbf{M} \otimes_{\mathbb{K}} \mathbb{S}(\mathbf{M})\) whose restriction to \(\mathbf{M}\) is the canonical homomorphism \(x \mapsto x \otimes 1\) ; every K-derivation \(\mathbf{D}: \mathbb{S}(\mathbf{M}) \to \mathbf{E}\) can therefore be written uniquely as \(h \circ \mathbf{D}_0\) with

\[
h \in \operatorname{Hom} _ {\mathbf {S} (\mathbf {M})} (\mathbf {M} \otimes_ {\mathbf {K}} \mathbf {S} (\mathbf {M}), \mathbf {E}).
\]

By the uniqueness of a solution of a universal mapping problem, it is seen that there exists a unique S(M)-module isomorphism

\[
\omega : \mathbf {M} \otimes_ {\mathbf {K}} \mathbf {S} (\mathbf {M}) \rightarrow \Omega_ {\mathbf {K}} (\mathbf {S} (\mathbf {M}))
\]

such that \(\mathbf{D}_0\circ \omega = d_{\mathbb{S}(\mathbf{M}) / \mathbb{K}}\) ; in other words, for all \(x\in \mathbf{M}\) \(\omega (x\otimes 1) = dx.\) In particular, if \(\mathbf{M}\) is a free K-module with basis \((e_{\lambda})_{\lambda \in \mathbf{L}},\Omega_{\mathbb{K}}(\mathbb{S}(\mathbf{M}))\) is a free S(M)-module with basis the set of differentials \(de_{\lambda}\) . Consider in particular the case where \(\mathrm{L} = [1,n]\) , so that \(\mathbb{S}(\mathbf{M})\) is identified with the polynomial algebra \(\mathbf{K}[\mathbf{X}_1,\dots ,\mathbf{X}_n]\) (§ 6, no. 6); for every polynomial \(\mathrm{P}\in \mathrm{K}[\mathrm{X}_1,\dots ,\mathrm{X}_n]\) , we can write uniquely

\[
d \mathrm{P} = \sum_ {i = 1} ^ {n} \mathrm{D} _ {i} \mathrm{P}. d \mathrm{X} _ {i}
\]

with \(D_{i}P \in K[X_{1}, \ldots, X_{n}]\) and, by virtue of the above, the mappings \(P \mapsto D_{i}P\) are the K-derivations of \(K[X_{1}, \ldots, X_{n}]\) corresponding to the coordinate forms on \(\Omega_{K}(S(M))\) for the basis \((dX_{i})\) ; we also write \(\frac{\partial P}{\partial X_{i}}\) instead of \(D_{i}P\) and this is called the partial derivative of P with respect to \(X_{i}\) .

